% ============================================================
%  Chapter fragment -- included by main.tex
%  (not a standalone document; compile main.tex to build the PDF)
% ============================================================
\begin{multicols*}{3}
\raggedcolumns

\sheettitle{11 -- Mobile Apps Security}{SWS2 -- ZHAW / Tellenbach}

% ==================== INTRO ====================
\section{Intro -- OWASP Mobile Top 10 \hl{(!)}}
Devs make the same security mistakes as in any app. Web has \kw{OWASP Top 10}; mobile has the \kw{OWASP Mobile Top 10} (M1--M10: improper platform usage, insecure data storage, insecure comm, insecure auth, insufficient crypto, insecure authz, client code quality, code tampering, reverse engineering, extraneous functionality).\\
It's \kw{less mature} (redundant categories) $\Rightarrow$ this chapter uses an \kw{adapted version}, \kw{6 categories}:
\begin{itemize}
\item \kw{Data Storage \& Leakage}, \kw{Secure Communication}, \kw{Auth \& Authorization}, \kw{Inter-App Communication}, \kw{Client-Side Injection}, \kw{Reverse Engineering}.
\end{itemize}
\textit{For each: what can go wrong + best practices.} \hl{Recurring theme: any in-app security measure can be circumvented by a determined attacker (esp. with root).}

% ============================================================
%  1. DATA STORAGE & LEAKAGE
% ============================================================
\section{1. Data Storage \& Leakage \hl{(!)}}
Storing sensitive data securely + preventing leakage. Leak channels: \kw{file system}, \kw{copy/paste buffer}, \kw{screenshots}, \kw{system log}, \kw{3rd-party} (analytics).
\subsection{Where to store data (3 options)}
{\footnotesize
\setlength{\tabcolsep}{3pt}
\renewcommand{\arraystretch}{1.1}
\arrayrulecolor{zhawblue}
\noindent\begin{tabular}{@{}>{\raggedright\arraybackslash}p{0.27\linewidth}>{\raggedright\arraybackslash}p{0.65\linewidth}@{}}
\rowcolor{zhawblue}
{\color{white}\bfseries Option} & {\color{white}\bfseries Notes}\\
\rowcolor{tableblue}
\kw{App-specific area} & \term{Documents} (iOS) / \term{shared\_prefs} (Android). Sandbox protects it; \hl{but root / jailbreak / backups expose it}.\\
\rowcolor{white}
\kw{Secure store} & iOS \kw{Keychain}/\kw{Secure Enclave}, Android \kw{Keystore}. Not in unencrypted backups; \hl{root still gets content}. Enclave best but \kw{public-key only}.\\
\rowcolor{tableblue}
\kw{Shared place} & Android \kw{MediaStore} / \term{/sdcard} (unprotected). \hl{Bad} unless data is meant to be shared.\\
\end{tabular}}
\textit{Android example: SharedPreferences write \term{user}/\term{password} into \term{...\_preferences.xml} -- other apps blocked by sandbox, but \kw{root} reads it \& it's in backups.}

\subsection{Storing sensitive data -- best practice \hl{(!)}}
\kw{Paranoid view} (\hl{store nothing sensitive, never store credentials}) is \kw{unrealistic}: users want \kw{convenience}, and \hl{forcing re-entry $\Rightarrow$ users pick weaker passwords}. So \kw{balance usability vs security}:
\begin{itemize}
\item \hl{Don't let users store passwords for highly critical services} (e-banking, crypto wallets, password managers).
\item For most apps it's \kw{reasonable} to store credentials in the \kw{app-specific FS} (sandbox makes access hard); use the \kw{Keychain/Keystore}; \kw{not} on \term{/sdcard}.
\item Extra encryption: \kw{hardcoded secret} $\to$ \hl{useless vs reverse engineering}; \kw{user-supplied secret at startup} $\to$ secure but \kw{not user-friendly}.
\end{itemize}
\kw{Be aware of the risks}: root malware reads all; \kw{unencrypted backups}; jailbreak/root raises risk; \hl{lending the phone} -- they use stored creds, or even \kw{root $\to$ extract creds $\to$ unroot} in minutes.

\subsection{Leakage channels}
\begin{itemize}
\item \kw{Copy/paste}: clipboard is \kw{global} (all apps read it) $\Rightarrow$ a Trojan reads copied \kw{passwords} (e.g. from a password manager). \term{Facebook app read clipboard to "offer" posting a copied link.}
\item \kw{Screenshots} -- \kw{by user}: stored in photo library, readable by apps with permission $\Rightarrow$ login/credit-card screen leaks. \kw{By OS}: snapshot when app \kw{backgrounded} (for app switcher) -- not app-readable but \kw{sits in FS}.
\item \kw{Logs}: any app can write; not app-readable but reachable via \kw{Xcode/adb}, root, jailbreak. \term{Banking app logged credentials -- forgot to disable for release.}
\end{itemize}
\subsection{Preventing leakage -- best practice}
\kw{Disable copy} on sensitive fields (both OS, programmatic); \kw{prevent OS screenshot} $\to$ blank image (both OS); \kw{block user screenshots} (\hl{Android only, impossible on iOS}); \kw{never log sensitive data} (turn off dev logging); \kw{check 3rd-party libs} (logging/analytics).

% ============================================================
%  2. SECURE COMMUNICATION
% ============================================================
\section{2. Secure Communication \hl{(!)}}
Apps talk to backends over \kw{HTTP} $\Rightarrow$ secure with \kw{TLS (HTTPS)}.
\begin{itemize}
\item \kw{Android 9+}: HTTPS \kw{enforced} (HTTP needs explicit opt-in). \kw{iOS}: HTTP needs exceptions in \kw{Info.plist} (since Jan 2017).
\item TLS is \kw{easy} -- supported by all languages/platforms.
\end{itemize}
\subsection{TLS best practices}
\kw{Disable insecure versions} (SSL, TLS 1.0) \& \kw{weak ciphers} (DES, RC4, MD5) -- usually \kw{secure by default} on recent runtimes. Use \kw{certs from established CAs}. \hl{Don't allow the user to override a failed cert check} (a failure may signal an attack).
\subsection{The CA problem $\to$ Certificate Pinning \hl{(!)}}
OS has a \kw{trust store} of many CAs. \hl{An attacker only needs ONE vulnerable CA} to get a valid cert for your server $\Rightarrow$ impersonation / \kw{MITM}. Since a mobile app talks to \kw{known servers only}, it can \kw{pin}:
{\footnotesize
\setlength{\tabcolsep}{3pt}
\renewcommand{\arraystretch}{1.1}
\arrayrulecolor{zhawblue}
\noindent\begin{tabular}{@{}>{\raggedright\arraybackslash}p{0.20\linewidth}>{\raggedright\arraybackslash}p{0.34\linewidth}>{\raggedright\arraybackslash}p{0.34\linewidth}@{}}
\rowcolor{zhawblue}
{\color{white}\bfseries Pin} & {\color{white}\bfseries Attacker must} & {\color{white}\bfseries Trade-off}\\
\rowcolor{tableblue}
\kw{CA cert} & get a fraudulent cert \kw{from the same CA} & survives server cert renewal; \kw{update app if CA changes} (rare)\\
\rowcolor{white}
\kw{Server cert} & get the \kw{server's private key} (or same-hash cert) & \kw{strongest}; \kw{update app on every cert change} (Let's Encrypt = 90d!)\\
\rowcolor{tableblue}
\kw{Own check} & self-signed server cert pinned in code & \kw{full control}, no external CA; rely on nothing external\\
\end{tabular}}
\textit{Pinning = hardcode the cert (or its hash) in the app, compare after TLS handshake.}

% ============================================================
%  3. AUTH & AUTHORIZATION
% ============================================================
\section{3. Authentication \& Authorization \hl{(!)}}
\begin{defbox}
\hl{Classic flaw:} app hides a screen behind login, but the \kw{video/data request to the server carries NO auth/authz} $\Rightarrow$ attacker requests files directly = free access. \kw{Hiding a UI screen $\ne$ protecting the function.} Cause: weak backend access control + believing UI = security.
\end{defbox}
\subsection{Weak auth on mobile}
Apps often \kw{relax} auth for convenience: \kw{shorter/weaker passwords}, \kw{2nd factor used less}. $\Rightarrow$ easier identity theft (guessing, offline cracking). \hl{Best practice: use the SAME auth strength as web}; require equally strong passwords \& \kw{store them} for usability (except highly sensitive, e.g. e-banking).
\subsection{Offline auth \& authz \hl{(!)}}
App checks credentials \kw{locally, no server} (e.g. game unlock code). \hl{Generally impossible to do securely} because the \kw{attacker has the whole app}:
\begin{itemize}
\item Content + check-logic + key must all be \kw{in the app}.
\item E.g. \kw{AC = HMAC\textsubscript{key}(username)}: attacker \kw{extracts the key} (make valid codes) \kw{OR} \kw{binary-patches} the check to always pass.
\item Anti-RE \kw{slows} but won't stop a committed attacker.
\end{itemize}
\subsection{Best practice}
\kw{Server-side access control}: after login, server issues an \kw{access token} sent on \kw{every request}; server checks identity + permission. \hl{Don't use offline auth/authz} unless the protected value is low. Don't allow weaker passwords; use a \kw{2nd factor} even if weaker than on a PC.
\subsection{2nd-factor options (vs mTAN)}
\kw{mTAN} = login code via SMS. Alternatives:
\begin{itemize}
\item \kw{Shared secret} app$\leftrightarrow$server (cheap): secret stored at registration, sent with login, server rotates it each login.
\item \kw{Separate authenticator app}: server sends a \kw{challenge} only the personalized auth app can solve (via inter-app comm, user confirms) $\to$ response forwarded to server.
\item \hl{Security of all $\approx$ mTAN}: to get the 2nd factor the attacker needs \kw{device access} (physical / malware) -- same as grabbing mTAN codes.
\end{itemize}

% ============================================================
%  4. INTER-APP COMMUNICATION
% ============================================================
\section{4. Inter-App Communication \hl{(!)}}
Both OS support it: \kw{URL schemes} (\term{zhaw://...}); Android also \kw{explicit Activity calls (intents)}. Main risks:
\begin{itemize}
\item Receiver \kw{processes data without validating} it (see client-side injection).
\item App \hl{exposes comm "accidentally"} $\to$ \kw{unlikely on iOS} (needs handling code) but \hl{likely on Android} (just a little \term{AndroidManifest.xml} config, no code).
\item \kw{Malicious app hijacks} the channel -- critical if sensitive data flows.
\end{itemize}
\subsection{Android specifics}
A component (Activity/Service/ContentProvider) is callable by other apps if:
\begin{itemize}
\item \kw{android:exported="true"} $\to$ callable by \kw{explicit intent} (package/class name, e.g. \term{...PostLogin}), \kw{OR}
\item it has an \kw{intent-filter} (defines URL scheme/action, e.g. \term{...VIEW\_CREDS}) $\to$ callable by \kw{implicit intent}.
\end{itemize}
\begin{defbox}
\hl{InsecureBank example:} a \kw{copy-paste mistake} marked the \kw{post-login} Activity \term{exported="true"} $\Rightarrow$ other apps / a user via \kw{adb} call it \kw{directly}, \kw{skipping login} = access to authenticated-only functions.
\end{defbox}
\subsection{Best practice}
\kw{Always validate} received data. Only set \kw{intent-filter / exported=true} if you really want the component public. \kw{Sensitive data}: a malicious app can register the \kw{same URL scheme} -- Android shows an \kw{app-chooser}, \hl{iOS gives it to the first registrant (since iOS 11)}; use \kw{explicit intents} to avoid the pop-up, and \kw{check the caller's identity} (App ID / package).

% ============================================================
%  5. CLIENT-SIDE INJECTION
% ============================================================
\section{5. Client-Side Injection \hl{(!)}}
Like server-side injection (SQLi, OS/XML) \kw{but inside the app, on the device}. Triggered \kw{by the user} or \kw{by another app} (inter-app). Two examples:
\subsection{Local SQL injection}
Apps use local \kw{SQLite} DBs. If the user controls part of a query $\to$ \kw{adapt its semantics}:
\begin{itemize}
\item \kw{Bypass local login}: \term{' OR 1=1 --}
\item \kw{Dump the DB}: \term{' or 1=1 --} in a search $\to$ returns all rows (all users, passwords, credit cards).
\end{itemize}
\hl{Risk usually smaller than server-side}: local DB often holds data \kw{the user already knows} \& is reachable anyway (adb). \kw{But} a threat when: used for \kw{local authentication} (lent phone), or the app \kw{accepts queries from other apps} (e.g. a health app) $\to$ another app uses SQLi to read \kw{arbitrary data}.
\subsection{Local File Inclusion}
App opens a file by a \kw{user-supplied name}; no validation $\Rightarrow$ \kw{path traversal} (\term{../shared\_prefs/...\_preferences.xml}) reads files outside the intended folder. Also via \kw{URLs}: app fetches a user URL, no validation $\to$ attacker uses a \kw{file://} URL (e.g. \term{file:///data/data/.../shared\_prefs/...}) to read local files. \hl{Mainly critical with inter-app communication} (normal user is sandbox-limited anyway).
\subsection{Best practice}
Same as server-side: \kw{prepared statements} (not string concatenation) for SQL; \kw{input validation} via \kw{whitelisting} (e.g. username = $\le$10 lowercase chars; filename = $\le$20 letters + one dot); for URLs \kw{accept only expected schemes} (http/https, \hl{not file}).

% ============================================================
%  6. REVERSE ENGINEERING
% ============================================================
\section{6. Reverse Engineering \hl{(!)}}
Goal: understand internals + \kw{extract hidden info}. Attacker motivation: find \kw{hardcoded secrets/keys}, \kw{hidden functionality}, \kw{locate \& defeat security functions} (offline access control), steal \kw{IP/algorithms}, clone the app (Trojan).
\subsection{Attacker's toolkit}
\begin{itemize}
\item \kw{strings} -- dump printable strings from a binary $\to$ spot passwords/URLs/commands (\term{olsdfgad;lh} "smells like a password").
\item \kw{Static analysis}: decompilers/disassemblers (Android Studio, JD-GUI, Hopper, IDA-Pro).
\item \kw{Dynamic analysis}: debuggers (gdb, jdb) -- watch runtime behaviour.
\item Then \kw{manipulate}: \kw{binary-patch} access control or jailbreak/root-detection to a \kw{no-op}; hook methods at runtime (\kw{cycript}/iOS, \kw{Frida}/Android).
\end{itemize}
\subsection{Best practice \hl{(!)}}
\hl{Disclaimer: anti-RE only slows the attacker, never stops a determined one.}
\begin{itemize}
\item \kw{Hardcoded secrets}: OK to \kw{encrypt local files "a bit"}; \hl{NOK to protect server admin functions} -- they'll be found.
\item \kw{Code obfuscation} (for valuable apps): harder to decompile. \hl{Esp. needed on Android} (Dalvik bytecode decompiles easily); ObjC/Swift harder (compiler optimization). Tools: \kw{ProGuard/DexGuard} (Android), \kw{iXGuard/PreEmptive} (iOS).
\item \kw{Anti-debugging}: iOS \term{SEC\_IS\_BEING\_DEBUGGED\_RETURN\_NIL()}, Android \term{Debug.isDebuggerConnected()} (gdb attach $\to$ segfault). Removable by \kw{binary patching}.
\item \kw{Jailbreak/root detection}: check for typical files (\term{/Applications/Cydia.app}, \term{/system/app/Superuser.apk}) $\Rightarrow$ raises the bar. \hl{But circumvented} by binary patching or hooking tools (\kw{Liberty Lite}/iOS, \kw{RootCloak}/Android) that fake a "clean" response.
\end{itemize}

% ==================== SUMMARY ====================
\section{Summary \hl{(!)}}
\begin{defbox}
Mobile-app flaws split in two: \kw{web-like} (secure comm, auth/authz, client-side injection -- fix as in web: TLS+pinning, server-side access control + tokens, prepared statements + whitelisting) and \kw{app-specific} (data storage/leakage, inter-app comm, reverse engineering). \hl{Golden rule: any in-app security measure can be circumvented by a determined attacker} -- offline auth, hardcoded-key encryption, and anti-RE only \kw{slow} attacks. $\Rightarrow$ keep secrets \& access control \kw{server-side}.
\end{defbox}

\end{multicols*}
